\documentclass[aip,jcp,preprint,longbibliography]{revtex4-2}

\usepackage{amsmath,amssymb,bm}
\usepackage{booktabs}
\usepackage{hyperref}
\hypersetup{
  colorlinks=true,
  citecolor=blue,
  linkcolor=blue,
  urlcolor=blue
}

\begin{document}

\title{Compressed Krylov Algorithm}

\author{Author Name}
\affiliation{Department of Chemistry, Institution Name, City, State ZIP, Country}

\date{\today}

\begin{abstract}
We present a draft formulation of a determinant-driven \emph{compressed Krylov algorithm} for strongly correlated electronic-structure calculations in active spaces. The central idea is to replace the exact Hamiltonian action required by a Krylov method with a selected-configuration projected matvec, compress each resulting vector to a bounded determinant support, and recover the approximate ground-state energy through a generalized Rayleigh-Ritz procedure. This construction positions the method between modern selected configuration interaction approaches, such as CIPSI, heat-bath CI, semistochastic heat-bath CI, and ASCI, and classical Krylov eigensolvers such as Lanczos and Davidson. In contrast to selected CI, which adapts a variational determinant list directly, the compressed Krylov algorithm adapts a small set of sparse Krylov vectors. In contrast to standard Krylov methods, the basis vectors are explicitly truncated and therefore are not assumed to remain orthonormal, which makes the generalized eigenvalue problem essential. The current implementation uses a matvec-first update: a connected single/double candidate pool is generated from the largest determinants in the current vector, projected amplitudes are evaluated against the full current support, determinants are ranked by projected-matvec or orthogonalized-residual amplitude rather than PT2 denominators, and compression is delayed until after orthogonalization. An alternate threshold-only compression mode is now also available, in which propagation and storage are controlled directly by coefficient thresholds. We summarize the algorithmic structure, its thick-restart and stochastic variants, and validated benchmarks from the current implementation.
\end{abstract}

\maketitle

\section{Introduction}
Accurate treatment of strong correlation in molecules continues to require methods that can approach full configuration interaction (FCI) quality within chemically relevant active spaces. Determinant-based selected configuration interaction (SCI) algorithms remain among the most successful routes because they exploit the sparsity structure of the wave function while retaining a clear variational picture. Canonical examples include CIPSI,\cite{huron1973cipsi} heat-bath configuration interaction (HCI),\cite{holmes2016hci} semistochastic heat-bath configuration interaction (SHCI),\cite{sharma2017shci} and adaptive sampling configuration interaction (ASCI).\cite{tubman2016asci,tubman2020asci}

These approaches are naturally complemented by determinant-driven software environments such as Quantum Package,\cite{garniron2019qp2} where variational selected spaces and perturbative corrections can be pushed to very large sizes. At the same time, alternative many-body solvers such as FCIQMC\cite{booth2009fciqmc} and DMRG\cite{white1992dmrg,chan2011dmrgchem} provide different paths to large-scale exact diagonalization.

The present work asks a different question: can one borrow the spectral acceleration of Krylov-subspace eigensolvers while preserving the determinant selectivity that makes SCI practical? Classical Krylov methods, from Lanczos\cite{lanczos1950} to Davidson,\cite{davidson1975} assume repeated access to a matrix-vector product. In a determinant basis, however, forming $H\bm q_k$ exactly may be substantially more expensive than constructing a carefully selected local approximation. The compressed Krylov algorithm (CKA) is designed around this observation.

The key distinction from SCI is conceptual. SCI adapts a variational determinant expansion and diagonalizes in that determinant space. CKA instead adapts a \emph{small number of sparse Krylov vectors}. Each vector is represented on a compressed determinant support, and each Hamiltonian action is approximated through a selected-space contraction that is designed to reproduce $H\bm q_k$ accurately on the retained support. Because the vectors are compressed after orthogonalization, they are not guaranteed to remain orthonormal, and the Ritz extraction must therefore solve a \emph{generalized} eigenvalue problem rather than an ordinary one.

\section{Method}
\subsection{Compressed Krylov basis}
Let $\{\bm q_k\}$ denote sparse determinant-space vectors intended to approximate a Krylov sequence for the electronic Hamiltonian $H$ in a fixed active-space basis. The starting vector may be the Hartree-Fock determinant or any precomputed multiconfigurational state. At iteration $k$, we define the Rayleigh quotient
\begin{equation}
E_k = \langle \bm q_k | H | \bm q_k \rangle.
\end{equation}
Instead of constructing $H\bm q_k$ exactly, the algorithm builds a selected determinant workspace $\mathcal{C}_k$ and computes a projected action
\begin{equation}
\widetilde{\bm v}_k \approx P_{\mathcal{C}_k} H \bm q_k,
\end{equation}
where $P_{\mathcal{C}_k}$ is the projector onto the selected determinant support.

In the current matvec-first implementation, $\mathcal{C}_k$ is assembled from three ingredients:
\begin{equation}
\mathcal{C}_k
=
\operatorname{supp}(\bm q_k)
\cup
\mathcal{E}_{\mathrm{SD}}(\mathcal{P}_k)
\cup
\mathcal{D}_k^{\mathrm{stoch}},
\end{equation}
where $\mathcal{P}_k$ is a capped parent set consisting of the largest-amplitude determinants in $\bm q_k$, $\mathcal{E}_{\mathrm{SD}}(\mathcal{P}_k)$ denotes their connected single and double excitations, and $\mathcal{D}_k^{\mathrm{stoch}}$ is an optional stochastic-discovery augmentation. The default deterministic production path uses only the first two terms. The projected amplitudes are then evaluated against the \emph{full current support} of $\bm q_k$,
\begin{equation}
\widetilde v_{k,I}
=
\sum_{J\in \operatorname{supp}(\bm q_k)}
H_{IJ} q_{k,J},
\label{eq:score}
\end{equation}
for $I \in \mathcal{C}_k$. This is the key change from a PT2-like SCI selector: the candidate builder is asked to approximate the Hamiltonian action itself, not to estimate a variational CI gain.

When previous basis vectors are available, the retained-support ranking is built from an Arnoldi-like residual,
\begin{equation}
\widetilde{\bm r}_k
=
\widetilde{\bm v}_k
-
\sum_{j \le k}
\bm q_j
\langle \bm q_j | \widetilde{\bm v}_k \rangle,
\end{equation}
and the default score is
\begin{equation}
\eta_I^{(k)} = |(\widetilde{\bm r}_k)_I|.
\end{equation}
If no previous basis is available, the fallback score is simply $|\widetilde v_{k,I}|$. The earlier PT2-style denominator score is retained only as a legacy ablation mode and is no longer the default propagation metric.

\subsection{Compression and generalized Ritz extraction}
The new implementation also delays compression. A larger temporary support
\begin{equation}
\widetilde \Omega_k = \operatorname{Top}_{n_{\mathrm{tmp}}}\!\left(\eta^{(k)}\right)
\end{equation}
is first selected, with $n_{\mathrm{tmp}} > n_{\mathrm{keep}}$ and with the current support optionally protected by a capped carry-over rule. The candidate update is built on $\widetilde \Omega_k$, orthogonalized against the previously stored sparse basis vectors, and only then compressed to a fixed support size $n_{\mathrm{keep}}$,
\begin{equation}
\bm q_{k+1}
 = \mathcal{N}\,
\mathcal{C}_{n_{\mathrm{keep}}}
\!\left[
\operatorname{orth}\!\left(P_{\widetilde \Omega_k}\widetilde{\bm v}_k;\{\bm q_j\}_{j\le k}\right)
\right],
\end{equation}
where $\mathcal{C}_{n_{\mathrm{keep}}}$ retains the largest-magnitude coefficients together with a capped subset of the previous support and $\mathcal{N}$ renormalizes the result. This change reduces the amount of irreversible truncation applied before the orthogonalization step and empirically improves parity with exact small-space reference calculations.

The current code also exposes a threshold-only compression policy. In that mode,
the propagation support is defined directly by
\begin{equation}
\mathcal{S}_k^{\mathrm{prop}}
=
\{J: |q_{k,J}| \ge \tau_{\mathrm{prop}}\},
\end{equation}
and the stored next vector is obtained from the orthogonalized residual by
\begin{equation}
\bm q_{k+1}
\propto
\mathcal{T}_{\tau_{\mathrm{store}}}\!\left(\widetilde{\bm r}_k\right),
\end{equation}
where $\mathcal{T}_{\tau_{\mathrm{store}}}$ keeps coefficients larger than a
user-selected storage threshold. If either threshold would otherwise remove the
entire vector, the largest-amplitude coefficient is retained as a safeguard so
that the iteration can continue.

If no compression were performed, a standard Krylov derivation would lead to an orthonormal basis and an ordinary projected eigenproblem. Compression breaks this property in general, even if an orthogonalization step is applied before truncation. The projected matrices are therefore
\begin{align}
(H_K)_{ij} &= \langle \bm q_i | H | \bm q_j \rangle, \\
(S_K)_{ij} &= \langle \bm q_i | \bm q_j \rangle,
\end{align}
and the Ritz problem is
\begin{equation}
H_K \bm c = E\, S_K \bm c.
\label{eq:genritz}
\end{equation}
Equation~\eqref{eq:genritz} is not a technical detail but a structural requirement of the compressed basis. In the current implementation, the overlap matrix is diagonalized, near-singular overlap modes are projected out, and the generalized Ritz problem is then solved in the retained overlap subspace. This stabilization became important once weakly compressed threshold-mode runs were used as a diagnostic limit.

\subsection{Selector backends and variants}
The compressed Krylov framework separates the basis update logic from the selected-space builder used to approximate $H\bm q_k$. Three selector backends are currently available:
\begin{enumerate}
  \item a deterministic connected-SD expansion mode, in which explicit single and double excitations are generated from a capped parent set and evaluated against the full current vector support,
  \item a stochastic discovery mode, which accumulates candidate support through repeated fixed-budget sampling rounds, and
  \item an explicit full-space fallback for validation.
\end{enumerate}

Separately, the selector score can be chosen from projected-matvec amplitudes, squared projected amplitudes, residual amplitudes, squared residual amplitudes, or the old PT2-like denominator score. The current default is the residual-amplitude form above.

This abstraction yields three practical solvers:
\begin{enumerate}
  \item a baseline CKA/CCIK propagation,
  \item a thick-restart variant that builds a local cycle of dimension $m_{\mathrm{cycle}}$, solves the Ritz problem, and recycles the lowest $n_{\mathrm{root}}$ Ritz vectors, and
  \item a stochastic variant in which candidate discovery is sampled rather than deterministically expanded.
\end{enumerate}

The thick-restart form is especially natural because compression already limits the support of each stored vector. Restarting after a modest local cycle allows the method to reuse a bounded number of sparse Ritz vectors while keeping the projected subspace manageable.

\section{Relationship to selected CI methods}
The compressed Krylov algorithm overlaps with the SCI family in its use of determinant ranking and selected projected contractions, but its adaptive object is different. Table~\ref{tab:positioning} summarizes this distinction.

\begin{table}[htbp]
\caption{Positioning of the compressed Krylov algorithm relative to common SCI families.}
\label{tab:positioning}
\begin{ruledtabular}
\begin{tabular}{p{0.22\linewidth} p{0.28\linewidth} p{0.23\linewidth} p{0.18\linewidth}}
Method family & Adaptive object & Typical ranking signal & Projected solve \\
\hline
CIPSI\cite{huron1973cipsi} & Variational determinant set & First-order or PT2 importance & CI diagonalization \\
HCI/SHCI\cite{holmes2016hci,sharma2017shci} & Variational determinant set & Heat-bath or semistochastic screening & CI diagonalization \\
ASCI\cite{tubman2016asci,tubman2020asci} & Variational determinant set & Coefficient-driven search and sorting & CI diagonalization \\
Compressed Krylov & Sparse Krylov vectors & Residual- or matvec-amplitude score, Eq.~\eqref{eq:score} & Generalized Ritz, Eq.~\eqref{eq:genritz}
\end{tabular}
\end{ruledtabular}
\end{table}

This viewpoint clarifies both the appeal and the risk of the method. The appeal is that one may obtain some of the spectral acceleration associated with Krylov iteration without abandoning determinant selection. The risk is that the method is \emph{not} an exact nested Krylov procedure after truncation. More projected matvecs do not automatically imply monotone improvement, because each step includes support restriction, orthogonalization against already compressed vectors, and a generalized projection onto a basis that is only approximately Krylov-like.

\section{Implementation outline}
The present draft reflects the \texttt{CompressK} codebase and its Quantum Package backend. The implementation uses sparse determinant vectors, selected projected contractions, and a generalized projected eigensolver over the compressed basis. In the current code, the default Quantum Package matvec backend is a determinant-driven \texttt{connected\_sd\_fullq} path: candidate determinants are generated by connected single/double expansion from a capped parent set, and projected amplitudes are evaluated against the full support of the current Krylov vector rather than against a truncated selector list. The same architecture supports:
\begin{itemize}
  \item a Python/PySCF path for dense active-space Hamiltonians,
  \item determinant-driven selection backends that mimic CIPSI-like expansion or stochastic discovery, and
  \item a Quantum Package plugin path for large determinant-driven production runs built on top of the CIPSI infrastructure in Quantum Package.\cite{garniron2019qp2}
\end{itemize}

The principal user parameters are the local Krylov dimension $m$ (or restart cycle size $m_{\mathrm{cycle}}$), the compressed support size $n_{\mathrm{keep}}$, the temporary support size $n_{\mathrm{tmp}}$, the parent cap used for connected-SD generation, and, for thick restart, the number of retained roots $n_{\mathrm{root}}$. In threshold mode, the user-facing compression controls reduce to the propagation threshold $\tau_{\mathrm{prop}}$ and storage threshold $\tau_{\mathrm{store}}$, while the local subspace dimensions remain the structural solver controls. The current codebase also includes rebuild-safe donor-backed Python-vs-QP parity scripts and cached exact-reference threshold-sweep drivers so that code changes can be validated without rerunning every benchmark from scratch.

\section{Validated benchmark snapshot}

\begin{table}[htbp]
\caption{Exact-Hamiltonian comparison for \(\mathrm{N}_2\), \(R=2.50\)~\AA, 6-31G, \((12\mathrm{o},12\mathrm{e})\), using the same donor-backed active-space FCIDUMP for all methods. The dense and thick CKA rows are the best completed weak-compression runs in the current Python implementation.}
\label{tab:n2dense}
\begin{ruledtabular}
\begin{tabular}{lccc}
Method & \(E_{\mathrm{tot}}\) (E$_\mathrm{h}$) & \(\Delta_{\mathrm{exact}}\) (mE$_\mathrm{h}$) & Size \\
\hline
Exact FCI & -108.7614 & 0.0000 & reference \\
DMRG & -108.7614 & 0.0000 & \(M_{\max}=3000\) \\
CIPSI & -108.7614 & 0.0008 & 32449 det \\
Dense CKA & -108.7614 & 0.0001 & 271420 uniq det \\
Thick CKA & -108.7614 & 0.0001 & 212729 uniq det \\
\end{tabular}
\end{ruledtabular}
\end{table}

Table~\ref{tab:n2dense} now shows that both dense and thick CKA can be driven essentially to the exact result on this donor-backed \(\mathrm{N}_2\) Hamiltonian. The best completed dense run used a full-candidate projected matvec, residual ranking, \(m=80\), \(n_{\mathrm{add}}=80000\), \(K_v=40000\), \(n_{\mathrm{keep}}=160000\), \(n_{\mathrm{gen}}=256\), and a temporary oversampling factor of \(2\). Its final total energy was \(-108.7614\) \(E_{\mathrm{h}}\), only \(0.0001\) m\(E_{\mathrm{h}}\) above exact FCI, with \(271420\) unique stored determinants across the full Krylov basis and \(12409320\) determinant entries when all \(80\) stored vectors are counted with multiplicity. The best completed thick run used \(m_{\mathrm{cycle}}=40\), \(n_{\mathrm{cycle}}=40\), \(n_{\mathrm{root}}=8\), the same \(n_{\mathrm{add}}\), \(K_v\), \(n_{\mathrm{keep}}\), and generator cap, and reached \(-108.7614\) \(E_{\mathrm{h}}\), also \(0.0001\) m\(E_{\mathrm{h}}\) above exact, with \(212729\) unique stored determinants and \(6562294\) determinant entries across the retained restart basis. The completed CIPSI comparison on the same Hamiltonian reached \(-108.7614\) \(E_{\mathrm{h}}\), \(0.0008\) m\(E_{\mathrm{h}}\) above exact, with \(32449\) saved determinants, so the best completed dense and thick CKA results are both lower variationally than the completed CIPSI result on this point. The same system is also diagnostically useful: a shorter dense run with \(m=15\) but much weaker compression \((n_{\mathrm{keep}}=120000)\) was worse by \(72.8322\) m\(E_{\mathrm{h}}\), which indicates that, for this case, Krylov-subspace dimension is more important than simply weakening compression at fixed \(m\).

\begin{table}[htbp]
\caption{Threshold-mode validation summary for the current implementation. Exact-reference dense/thick gates are reported for \(\mathrm{N}_2\), 6-31G, \((10\mathrm{o},10\mathrm{e})\), \(R=2.85\)~\AA. Donor-backed Python-vs-QP parity is reported for the current \texttt{connected\_sd\_fullq} plugin backend.}
\label{tab:thresholdgate}
\begin{ruledtabular}
\begin{tabular}{p{0.32\linewidth}p{0.28\linewidth}p{0.28\linewidth}}
Check & System / setting & Result \\
\hline
Dense zero-threshold exact gate & \(\mathrm{N}_2\), \((10\mathrm{o},10\mathrm{e})\) & max error \(= 5.05 \times 10^{-5}\) m\(E_{\mathrm{h}}\) \\
Thick zero-threshold exact gate & \(\mathrm{N}_2\), \((10\mathrm{o},10\mathrm{e})\) & max error \(= 9.87 \times 10^{-4}\) m\(E_{\mathrm{h}}\) \\
Zero-threshold QP parity & \(\mathrm{N}_2\), \((10\mathrm{o},10\mathrm{e})\) & \(q_{\mathrm{next}}\) L\(_2 = 7.95 \times 10^{-15}\), energy gap \(= 1.11 \times 10^{-5}\) m\(E_{\mathrm{h}}\) \\
Zero-threshold QP parity & \(\mathrm{N}_2\), \((12\mathrm{o},12\mathrm{e})\) & \(q_{\mathrm{next}}\) L\(_2 = 1.86 \times 10^{-16}\), energy gap \(= 4.88 \times 10^{-6}\) m\(E_{\mathrm{h}}\) \\
\end{tabular}
\end{ruledtabular}
\end{table}

Table~\ref{tab:thresholdgate} summarizes the current threshold-mode certification surface. On the Python side, dense and thick zero-threshold runs remain within chemical accuracy of exact FCI on a hard \(\mathrm{N}_2\) point. On the QP side, donor-backed parity with the rebuilt \texttt{connected\_sd\_fullq} plugin remains near machine precision for the propagated next vector on both \((10\mathrm{o},10\mathrm{e})\) and \((12\mathrm{o},12\mathrm{e})\) active spaces. In addition, the threshold sweep on \(\mathrm{N}_2\), 6-31G, \((8\mathrm{o},10\mathrm{e})\), \(R=2.50\)~\AA\ shows a clear algorithmic pattern: weakening the storage threshold produces only gradual degradation, whereas tightening the propagation threshold eventually degrades the energy much more rapidly because it removes determinants before the projected Hamiltonian action is assembled. The best weak-threshold points in the current sweep gave dense and thick errors of \(4.65 \times 10^{-6}\) and \(4.59 \times 10^{-7}\) m\(E_{\mathrm{h}}\), respectively, for the storage-threshold scan, while the best propagation-threshold points were \(9.64 \times 10^{-6}\) and \(5.44 \times 10^{-7}\) m\(E_{\mathrm{h}}\). Thick restart is therefore slightly more robust across the sweep, but both modes recover the exact limit when the thresholds are sufficiently weak.

The large fixed-orbital Cr\(_2\) production campaign is intentionally not included in the present draft benchmark tables, because that comparison has not yet been rerun end-to-end against the fully current code path after the finalized threshold-mode and generalized-Ritz hardening work. The current manuscript therefore restricts itself to validated \(\mathrm{N}_2\) benchmarks that are already fully reproducible with the maintained code.

\section{Discussion}
Three observations are worth emphasizing.

First, the generalized Ritz step is mandatory whenever the basis vectors are truncated. Any implementation that treats the compressed basis as exactly orthonormal is, at best, making an uncontrolled approximation. The current code explicitly constructs both $H_K$ and $S_K$, filters near-singular overlap modes, and then solves Eq.~\eqref{eq:genritz} in the retained overlap subspace.

Second, the selected-space projected matvec is the central algorithmic approximation. The new implementation makes this statement literal: candidate generation is now organized around reproducing $H\bm q_k$ on a connected-SD support, and the default ranking signal is the orthogonalized residual amplitude rather than a PT2 denominator heuristic. Better candidate builders and residual selectors should therefore improve the quality of each projected $H\bm q_k$ and, by extension, the quality of the compressed Krylov space.

Third, thick restart is natural but not automatically sufficient. Because the method is no longer an exact Lanczos recurrence after compression, restart quality depends on which Ritz information is retained and how the projected supports are rebuilt at the next cycle. The threshold-mode sweeps also show that the accuracy bottleneck is more sensitive to propagation truncation than to storage truncation, which is a useful design cue for future restart and support-management strategies.

\section{Conclusions}
The compressed Krylov algorithm provides a simple organizing principle for combining determinant selection with Krylov-subspace acceleration. Its defining features are:
\begin{enumerate}
  \item selected-space approximations to the Hamiltonian action,
  \item explicit compression of each basis vector to bounded determinant support, and
  \item generalized Rayleigh-Ritz extraction in the resulting nonorthogonal basis.
\end{enumerate}

In that sense, the method stands between selected CI and classical Krylov iteration rather than replacing either one outright. The current implementation already supports deterministic, thick-restart, stochastic, and threshold-only variants, together with donor-backed Python-vs-QP parity checks and exact-reference threshold sweeps. The next stage of the project is a broader benchmark campaign against high-budget CIPSI, SHCI, and ASCI calculations, together with more careful study of restart strategies, threshold policies, and orbital optimization within the compressed-Krylov picture.

\bibliography{compressed_krylov_algorithm}

\end{document}
